\chapter{آزمایش‌ها}
\label{ch:experiments}

این فصل چیدمان تجربی و طراحی آزمایش‌های \pnum{1} تا \pnum{5} را شرح می‌دهد.

%==============================================================================
\section{چیدمان تجربی}
%==============================================================================

\subsection{پیاده‌سازی و معماری نرم‌افزار}

پیاده‌سازی آزمایش‌ها با ابزارها و کتابخانه‌های زیر انجام شده است:

\begin{itemize}
\item \textbf{زبان و هسته}: \lr{Python 3.10+}.
\item \textbf{یادگیری عمیق}: \lr{PyTorch} برای رمزگذار \lr{ETHOS} و آموزش کم‌نمونه.
\item \textbf{یادگیری ماشین کلاسیک}: \lr{scikit-learn} (\lr{Logistic Regression}، \lr{Random Forest}، معیارها)، \lr{XGBoost}، \lr{LightGBM}، \lr{CatBoost}.
\item \textbf{داده}: \lr{NumPy}، \lr{Pandas} برای پیش‌پردازش و مهندسی ویژگی.
\item \textbf{فدرال}: پیاده‌سازی سفارشی \lr{FedAvg} (طراحی سازگار با \lr{PySyft} برای گسترش آینده).
\item \textbf{تصویرسازی}: \lr{Matplotlib}، \lr{Seaborn} برای منحنی‌های \lr{ROC} و منحنی‌های یادگیری.
\end{itemize}

آزمایش‌ها روی ایستگاه کاری استاندارد با \lr{CPU} چندهسته‌ای اجرا شدند. همه‌ی بذرهای تصادفی ثابت (\lr{seed=42}) برای بازتولیدپذیری تنظیم شده‌اند. برای نیازمندی‌ها و جزئیات محیط، پیوست~\ref{app:repro} را ببینید.

\subsection{پیکربندی داده}

\begin{itemize}
\item مجموعه‌داده: \lr{MIMIC\text{-}IV v3.1}.
\item اندازۀ کوهورت: \pnum{2,000} بیمار.
\item تقسیم‌ها: \pnum{64}\% آموزش (\pnum{1,280})، \pnum{16}\% اعتبارسنجی (\pnum{320})، \pnum{20}\% آزمون (\pnum{400}).
\item گره‌های بیماری: \pnum{5} (سپسیس، قلبی، تنفسی، کلیوی، دیابت).
\item ویژگی‌ها: \pnum{25} بعد (\pnum{22} پایه + \pnum{3} مهندسی‌شده).
\end{itemize}

\subsection{\texorpdfstring{پیکربندی \lr{ETHOS} و یادگیری کم‌نمونه}{پیکربندی ETHOS و یادگیری کم‌نمونه}}

\begin{itemize}
\item $K = 5$ پشتیبان، $Q = 15$ پرس‌وجو به‌ازای هر کلاس.
\item \pnum{500} اپیزود آموزشی به‌ازای هر دوره، \pnum{100} اپیزود اعتبارسنجی.
\item \pnum{50} دوره، صبر توقف زودهنگام \pnum{10}.
\item دما $\tau = 10$، زیان کانونی $\gamma = 2$.
\item نرخ یادگیری \pnum{0.001}، بهینه‌ساز \lr{Adam}.
\end{itemize}

\subsection{پیکربندی خط پایۀ کلاسیک}

\begin{itemize}
\item \lr{RandomizedSearchCV} با \prange{50}{80} آزمون (\lr{XGBoost}، \lr{Random Forest})، \pnum{50} آزمون (\lr{LightGBM}، \lr{CatBoost}).
\item اعتبارسنجی متقاطع لایه‌ای \pnum{5}-تایی روی آموزش+اعتبارسنجی.
\item امتیازدهی: \lr{AUROC}.
\item تنظیم آستانه روی اعتبارسنجی: $\theta \in [0.3, 0.7]$، گام \pnum{0.01}، بیشینه‌سازی \lr{F1}.
\end{itemize}

%==============================================================================
\section{فرضیات آزمایش}
%==============================================================================

هر آزمایش برای آزمون فرضیات مشخص طراحی شده است:

\begin{table}[htbp]
\centering
\caption{فرضیات آزمایش و پیامدهای مورد انتظار}
\footnotesize
\setlength{\tabcolsep}{2pt}
\renewcommand{\arraystretch}{1.14}
\begin{tabularx}{\linewidth}{@{}C{1.15cm}R{4.15cm}Y@{}}
\toprule
\textbf{\lr{Exp}} & \textbf{فرضیه} & \textbf{پیامد مورد انتظار} \\
\midrule
\pnum{1} & کم‌نمونه می‌تواند خطوط پایۀ کلاسیک را جبران کند & مقایسۀ \lr{AUROC} بین مدل‌ها \\
\pnum{2} & $K$ بزرگ‌تر عملکرد کم‌نمونه را بهبود می‌دهد & افزایش \lr{AUROC} با $K$ \\
\pnum{3} & فدرال هزینۀ ناچیز در برابر متمرکز دارد & شکاف کمتر از \pnum{2}\% \lr{AUROC} \\
\pnum{4} & وزن‌دهی شدت بحرانی است & حذف، \lr{AUROC} را به‌طور محسوسی کاهش می‌دهد \\
\pnum{5} & مدل به گره‌های بیماری دیده‌نشده تعمیم می‌یابد & \mbox{\lr{AUROC $>$ 0.55}} روی گره‌های نگه‌داشته‌شده \\
\bottomrule
\end{tabularx}
\label{tab:exp_hypotheses}
\end{table}

%==============================================================================
\section{توان آماری و اندازۀ نمونه}
%==============================================================================

با \pnum{400} نمونۀ آزمون و تقسیم کلاس \pnum{42}/\pnum{58}، تقریباً \pnum{168} مثبت و \pnum{232} منفی داریم. برای تشخیص اختلاف \lr{AUROC} به میزان \pnum{0.10} (مثلاً \pnum{0.76} در برابر \pnum{0.66}) با توان \pnum{80}\% و $\alpha = 0.05$، آزمون \lr{DeLong} تقریباً \prange{100}{200} نمونه به‌ازای هر گروه لازم دارد؛ مجموعۀ آزمون ما این شرط را برآورده می‌کند. بازه‌های اطمینان بوت‌استرپ (\pnum{1,000} نمونه‌برداری مجدد) برآوردهای پایدار می‌دهند؛ برای معیارهای کلیدی بازه‌ی اطمینان \pnum{95}\% گزارش می‌کنیم.

%==============================================================================
\section{مشخصات سخت‌افزار و نرم‌افزار}
%==============================================================================

آزمایش‌ها روی ایستگاه کاری استاندارد انجام شد: \lr{Intel Core i7} یا معادل، \pnum{16} تا \pnum{32} گیگابایت \lr{RAM}، \lr{GPU} انویدیا اختیاری (\lr{CUDA 11.x}). زمان‌های آموزش: \lr{Random Forest} حدود یک دقیقه؛ آموزش \lr{ETHOS} با \pnum{50} دوره حدود دو تا سه ساعت روی \lr{CPU} یا حدود \pnum{15} دقیقه با \lr{GPU}. نرم‌افزار: \lr{Python 3.10}، \lr{PyTorch 2.x}، \lr{scikit-learn 1.3+}، \lr{XGBoost 2.x}، \lr{pandas 2.x}. همه‌ی وابستگی‌ها در \texttt{\lr{requirements.txt}} فهرست شده‌اند (پیوست~\ref{app:repro}).

%==============================================================================
\section{فلسفۀ طراحی آزمایش}
%==============================================================================

آزمایش‌ها به‌صورت سلسله‌مراتبی طراحی شده‌اند: (\pnum{1}) ابتدا خطوط پایۀ قوی با \lr{ML} کلاسیک ساخته می‌شوند؛ (\pnum{2}) سپس رویکردهای عصبی و کم‌نمونه با آن‌ها مقایسه می‌شوند؛ (\pnum{3}) با حذف تدریجی اجزا، سهم هر بخش بررسی می‌شود؛ و (\pnum{4}) تعمیم فدرال و بین‌بیماری ارزیابی می‌گردد. بذرهای تصادفی و تقسیم‌های داده ثابت نگه داشته شده‌اند تا بازتولیدپذیری حفظ شود. هر آزمایش با پیکربندی مشخص اجرا و برآورد نقطه‌ای آن گزارش می‌شود؛ بازه‌های اطمینان بوت‌استرپ و آزمون‌های آماری در فصل~\ref{ch:results} آمده‌اند.

%==============================================================================
\section{رویۀ اجرای آزمایش}
%==============================================================================

همه‌ی آزمایش‌ها از رویۀ استاندارد پیروی می‌کنند: (\pnum{1}) بارگذاری دادۀ پیش‌پردازش‌شده (تقسیم‌های آموزش/اعتبارسنجی/آزمون)؛ (\pnum{2}) برای مدل‌های کلاسیک: برازش روی آموزش، تنظیم آستانه روی اعتبارسنجی، ارزیابی روی آزمون؛ (\pnum{3}) برای چارچوب \lr{ETHOS} و یادگیری کم‌نمونه: آموزش اپیزودیک روی آموزش، اپیزودهای اعتبارسنجی روی اعتبارسنجی، ارزیابی روی آزمون؛ (\pnum{4}) برای فدرال: پارتیشن بر اساس گرۀ بیماری، اجرای \lr{FedAvg} برای \pnum{20} دور، ارزیابی مدل سراسری روی آزمون؛ (\pnum{5}) ثبت معیارها (\lr{AUROC}، \lr{F1}، و غیره) و ذخیرۀ مدل‌ها. از دادۀ آزمون برای انتخاب مدل استفاده نمی‌شود. همه‌ی آزمایش‌ها از همان تقسیم‌های ثابت (\lr{seed=42}) استفاده می‌کنند.

%==============================================================================
\section{پیکربندی‌های مدل خط پایه}
%==============================================================================

پیکربندی تفصیلی برای هر خط پایه ارائه می‌شود؛ مقادیر فنی برای جلوگیری از شکستن جهت متن و برش حاشیه در یک بلوک لاتینِ چپ‌به‌راست دست‌نخورده می‌آید (با قلم \lr{monospace}). شماره‌های گزارشی بهترین‌ها (\pnum{200}، \pnum{28}، و امثال آن) در توضیح فارسی حفظ شده‌اند.
\begin{latin}\footnotesize\setlength{\parskip}{0.35em}\setlength{\parindent}{0pt}\ttfamily\sloppy\raggedright
\textbf{Logistic Regression.}\quad L2 penalty (C{=}1.0).\quad class\_weight{=}balanced.\quad max\_iter{=}1000.\\[-0.1em]
\textbf{Random Forest.}\quad Selected via RandomizedSearchCV:\quad n\_estimators{=}200;\quad max\_depth{=}28;\quad min\_samples\_split{=}5;\quad min\_samples\_leaf{=}2;\quad max\_features{=}log2.\\[-0.1em]
\textbf{XGBoost.}\quad n\_estimators{=}200;\quad max\_depth{=}6;\quad learning\_rate{=}0{.}05;\quad subsample{=}0{.}8;\quad colsample\_bytree{=}0{.}8.\\[-0.1em]
\textbf{LightGBM.}\quad n\_estimators{=}200;\quad max\_depth{=}6;\quad learning\_rate{=}0{.}05;\quad num\_leaves{=}31.\\[-0.1em]
\textbf{CatBoost.}\quad iterations{=}200;\quad depth{=}6;\quad learning\_rate{=}0{.}05.
\end{latin}
همه‌ی مدل‌های کلاسیک از همان بردار ویژگی جدولی \pnum{25}-بعدی و تقسیم‌های یکسان آموزش/اعتبارسنجی/آزمون استفاده می‌کنند. تنظیم آستانه روی مجموعۀ اعتبارسنجی برای همه‌ی مدل‌ها برای بیشینه‌سازی \lr{F1} انجام می‌شود؛ آستانۀ تنظیم‌شده برای ارزیابی نهایی روی آزمون اعمال می‌شود.

%==============================================================================
\section{معیارهای ارزیابی و گزارش}
%==============================================================================

برای همه‌ی مدل‌ها \lr{AUROC} (اصلی)، امتیاز \lr{F1} کلان، \lr{Precision}، \lr{Recall}، \lr{Accuracy} و \lr{MCC} گزارش می‌کنیم. \lr{AUROC} نسبت به آستانۀ تصمیم مستقل است و برای دادۀ نامتعادل مناسب‌تر از اتکا به یک آستانۀ ثابت است. امتیاز \lr{F1} کلان میانگین \lr{F1} به‌ازای هر کلاس را می‌گیرد تا سلطۀ کلاس اکثریت را کاهش دهد. همچنین بازه‌های اطمینان \pnum{95} درصد بوت‌استرپ (\pnum{1,000} نمونه‌برداری مجدد) برای معیارهای کلیدی گزارش می‌کنیم. آزمون‌های جفتی مانند \lr{DeLong} فقط زمانی معتبرند که بردارهای پیش‌بینی هم‌تراز روی همان نمونه‌های آزمون موجود باشد؛ بنابراین مقایسه‌های بین‌پروتکلی در این پایان‌نامه عمدتاً توصیفی تفسیر می‌شوند. همه‌ی معیارها روی مجموعۀ آزمون نگه‌داشته‌شده محاسبه می‌شوند و از دادۀ آزمون برای انتخاب مدل یا تنظیم ابرپارامتر استفاده نمی‌شود. برای ارزیابی کم‌نمونه، \pnum{100} اپیزود نمونه‌برداری می‌کنیم؛ هر اپیزود \lr{K} پشتیبان و \lr{Q} پرس‌وجو به‌ازای هر کلاس دارد. میانگین و انحراف معیار بین اپیزودها را گزارش می‌کنیم. برای مدل‌های کلاسیک یک‌بار روی کل مجموعۀ آزمون ارزیابی می‌کنیم. تفاوت پروتکل‌ها بدین معناست که معیارهای آزمایش \lr{K}-شات (در این مخزن حدود \prange{0.557}{0.564} \lr{AUROC}) مستقیماً با ارزیابی ثابت \lr{Exp1} برای \texttt{\lr{Proposed\_FewShot}} (\pnum{0.611}) قابل‌تعویض نیستند.

%==============================================================================
\section{آزمایش \pnum{1}: مقایسۀ مدل اصلی}
%==============================================================================

\subsection{هدف}

مقایسۀ چارچوب \lr{ETHOS} و مدل پیشنهادی کم‌نمونه با خطوط پایۀ کلاسیک و عصبی روی مجموعۀ آزمون.

\subsection{مدل‌ها}

\begin{itemize}
\item \texttt{\lr{LogisticRegression}}
\item \texttt{\lr{RandomForest}}
\item \texttt{\lr{XGBoost}}
\item \texttt{\lr{StandardNN}} (\lr{MLP}، بدون کم‌نمونه)
\item \texttt{\lr{FederatedNN\_NoFewShot}} (\lr{MLP} فدرال، بدون کم‌نمونه)
\item مدل پیشنهادی کم‌نمونه (\lr{ETHOS} + سر نمونه‌اولیه)
\end{itemize}

\subsection{خلاصۀ نتایج}

مدل پیشنهادی کم‌نمونه به \lr{AUROC} \pnum{0.611}، \lr{F1} \pnum{0.572}، و دقت \pnum{0.583} می‌رسد. بهترین مدل کلاسیک در همین پروتکل \lr{Random Forest} است که \lr{AUROC} \pnum{0.762}، \lr{F1} \pnum{0.679}، و دقت \pnum{0.718} دارد. این اعداد با جدول~\ref{tab:exp1} در فصل~\ref{ch:results} همگام‌اند.

%==============================================================================
\section{آزمایش \pnum{2}: حساسیت \lr{K}-شات}
%==============================================================================

\subsection{هدف}

هدف این آزمایش، سنجش اثر تعداد نمونه‌های پشتیبان \lr{K} بر عملکرد کم‌نمونه است. در این چارچوب، \lr{K} تعداد نمونۀ برچسب‌خورده برای هر کلاس در مجموعۀ پشتیبان است. افزایش \lr{K} اطلاعات بیشتری برای ساخت نمونه‌های اولیه فراهم می‌کند، اما هم‌زمان مسئله را از حالت کم‌نمونه دورتر می‌سازد. بنابراین \lr{K} را تغییر می‌دهیم تا این بده‌بستان بهتر روشن شود.

\subsection{جزئیات طراحی}

برای هر مقدار \lr{K} از میان \pnum{1}، \pnum{3}، \pnum{5}، \pnum{10} و \pnum{15}، اپیزودها با \lr{K} پشتیبان و \lr{Q=15} پرس‌وجو به‌ازای هر کلاس نمونه‌برداری می‌شوند. هر پیکربندی روی \pnum{100} اپیزود آزمون ارزیابی می‌شود و میانگین \lr{AUROC} و انحراف معیار گزارش می‌گردد. مدل اصلی \lr{ETHOS} و یادگیری کم‌نمونه برای آموزش از \lr{K=5} استفاده می‌کند؛ این آزمایش \lr{K} را فقط در زمان ارزیابی متغیر می‌کند تا حساسیت سنجیده شود. توجه: پروتکل ارزیابی اپیزودیک با ارزیابی اصلی \lr{Exp1} متفاوت است؛ از این رو مقادیر مطلق \lr{AUROC} در این آزمایش در بازۀ \prange{0.557}{0.564} می‌مانند و از عدد منجمد \pnum{0.611} پایین‌ترند.

\subsection{طراحی}

$K \in \{1, 3, 5, 10, 15\}$ را متغیر می‌کنیم و \lr{AUROC} و \lr{F1} را روی مجموعۀ آزمون اندازه می‌گیریم. هر پیکربندی روی چند اپیزود ارزیابی می‌شود.

\subsection{خلاصۀ نتایج}

\lr{AUROC} در بازۀ باریک \prange{0.557}{0.564} نوسان می‌کند. بهترین میانگین مشاهده‌شده \pnum{0.564} در \lr{K=3} است و \lr{K=1} نیز تقریباً هم‌تراز باقی می‌ماند؛ بنابراین در این جاروب اپیزودیک روند صعودی یکنواخت دیده نمی‌شود. جدول~\ref{tab:exp2} را ببینید.

%==============================================================================
\section{آزمایش \pnum{3}: فدرال در برابر متمرکز}
%==============================================================================

\subsection{هدف}

هدف این آزمایش، اندازه‌گیری هزینۀ عملکردی یادگیری فدرال است. در یادگیری فدرال، مدل بدون تجمیع نمونه‌های خام روی دادۀ توزیع‌شده آموزش می‌بیند؛ بنابراین پرسش اصلی این است که این قید تا چه اندازه به عملکرد آسیب می‌زند.

\subsection{طراحی}

چارچوب \lr{ETHOS} و یادگیری کم‌نمونه در دو چیدمان آموزش داده می‌شود: (\pnum{1}) متمرکز، که در آن همۀ داده‌ها در یک محل در دسترس‌اند و آموزش اپیزودیک استاندارد انجام می‌شود؛ و (\pnum{2}) فدرال، که در آن داده بر اساس بیماری به \pnum{5} گره تقسیم می‌شود و \lr{FedAvg} در \pnum{20} دور با \pnum{3} دورۀ محلی در هر دور اجرا می‌گردد. هر دو چیدمان از معماری، ابرپارامترها و تقسیم‌های داده یکسان استفاده می‌کنند. در حالت فدرال، هر گره مدل را محلی آموزش می‌دهد و به‌روزرسانی‌ها را به سرور می‌فرستد؛ سرور نیز آن‌ها را متناسب با اندازۀ گره تجمیع می‌کند. \lr{AUROC} و \lr{F1} روی مجموعۀ آزمون نگه‌داشته‌شده مقایسه می‌شوند. مجموعۀ آزمون سراسری و بدون پارتیشن است تا مقایسه منصفانه باقی بماند.

\subsection{خلاصۀ نتایج}

در خروجی نهایی این آزمایش، چیدمان متمرکز \lr{AUROC} \pnum{0.552} و \lr{F1} \pnum{0.543} دارد، در حالی که چیدمان فدرال \lr{AUROC} \pnum{0.547} و \lr{F1} \pnum{0.540} ثبت می‌کند؛ بنابراین افت مطلق حدود \pnum{0.005} \lr{AUROC} است. این نتیجه از نظر کاربردی کوچک است، اما به‌معنای برتری فدرال نیست؛ تنها نشان می‌دهد که در این شبیه‌سازی هزینه‌ی حفظ حریم خصوصی محدود مانده است. جدول~\ref{tab:exp3} را ببینید.

%==============================================================================
\section{آزمایش \pnum{4}: مطالعۀ حذف تدریجی}
%==============================================================================

\subsection{هدف}

هدف این آزمایش، شناسایی اجزایی است که بیشترین سهم را در عملکرد مدل دارند. در حذف تدریجی، هر بار یک جزء حذف می‌شود و تغییر عملکرد اندازه‌گیری می‌گردد. این کار کمک می‌کند مشخص شود کدام انتخاب‌های طراحی واقعاً حیاتی‌اند و کدام‌ها اثر محدودتری دارند.

\subsection{گونه‌ها}

\begin{itemize}
\item \texttt{\lr{full\_model}}: چارچوب کامل \lr{ETHOS} و یادگیری کم‌نمونه.
\item \texttt{\lr{no\_symptom\_tokens}}: جایگزینی با \lr{Logistic Regression} روی \pnum{25} ویژگی خام.
\item \texttt{\lr{no\_clinical\_bias}}: حذف سوگیری بالینی در بردارهای نهفته.
\item \texttt{\lr{no\_intensity\_weighting}}: حذف وزن‌دهی شدت.
\item \texttt{\lr{no\_temporal\_decay}}: حذف واگرایی زمانی.
\item \texttt{\lr{no\_few\_shot}}: \lr{NN} استاندارد (بدون آموزش اپیزودیک).
\end{itemize}

\subsection{خلاصۀ نتایج}

\texttt{\lr{full\_model}}: \lr{AUROC} \pnum{0.583}. \texttt{\lr{no\_symptom\_tokens}}: \pnum{0.724} (\lr{LR} روی ویژگی‌های خام). \texttt{\lr{no\_intensity\_weighting}}: \pnum{0.577}. در این جدول، حذف توکن‌های علائم ظاهراً عملکرد را بهتر می‌کند؛ بنابراین این حذف تدریجی بیشتر به محدودیت بازنمایی فعلی اشاره می‌کند تا تأیید بی‌قیدوشرط اجزای پیشنهادی. جدول~\ref{tab:exp4} را ببینید.

%==============================================================================
\section{آزمایش \pnum{5}: تعمیم بین‌بیماری}
%==============================================================================

\subsection{هدف}

ارزیابی تعمیم به گروه‌های بیماری دیده‌نشده. در چیدمان‌های فدرال یا چندسایتی، مدلی که روی برخی گروه‌های بیماری آموزش دیده ممکن است با دادۀ محلی محدود به دیگران تعمیم یابد.

\subsection{طراحی}

روی گره‌های \prange{1}{3} (سپسیس، قلبی، تنفسی) آموزش؛ آزمون روی گره‌های \pnum{4} و \pnum{5} (کلیوی، دیابت). این سناریویی را شبیه‌سازی می‌کند که بیمارستان یا گروه بیماری جدیدی به فدراسیون می‌پیوندد بدون دادۀ آموزش محلی. \texttt{\lr{Proposed\_FewShot}}، \lr{Logistic Regression}، \lr{XGBoost}، و \lr{Standard NN} را مقایسه می‌کنیم. یادگیری کم‌نمونه برای چنین سناریوهایی طراحی شده است؛ آزمون می‌کنیم آیا وقتی توزیع آزمون از توزیع آموزش متفاوت است بر مدل‌های کلاسیک برتری دارد.

\subsection{خلاصۀ نتایج}

پیشنهادی: \pnum{0.602} (\texttt{\lr{node\_4}}) و \pnum{0.478} (\texttt{\lr{node\_5}}). \lr{LR}: \pnum{0.668} و \pnum{0.658}. \lr{XGBoost}: \pnum{0.704} و \pnum{0.700}. در این اسنپ‌شات مدل‌های کلاسیک روی هر دو گره بهتر تعمیم می‌یابند و به‌ویژه روی گره‌ی دیابت افت مدل پیشنهادی چشمگیر است. جدول~\ref{tab:exp5} را ببینید.

%==============================================================================
\section{رویۀ جستجوی ابرپارامتر}
%==============================================================================

برای مدل‌های کلاسیک از \lr{RandomizedSearchCV} با \prange{50}{80} آزمون استفاده می‌کنیم. فضاهای جستجو در بلوک لاتینِ زیر به‌صورت یک‌خطیِ قابل‌شکستن آمده تا ارقام و پرانتز با متن فارسی قاطی نشوند؛ معیار امتیاز \lr{AUROC} است؛ اعتبارسنجی متقاطع لایه‌ای \pnum{5}‌تایی روی آموزش+اعتبارسنجی انجام می‌شود؛ بهترین پیکربندی با میانگین \lr{AUROC} روی چین‌ها انتخاب می‌شود؛ سپس آستانۀ تصمیم روی اعتبارسنجی برای بیشینه‌سازی \lr{F1} تنظیم و معیارهای نهایی روی آزمون گزارش می‌شوند. برای شبکه‌ی کامل پیوست~\ref{app:hyperparams} را ببینید.

\begin{latin}\footnotesize\setlength{\parskip}{0.35em}\setlength{\parindent}{0pt}\ttfamily\sloppy\raggedright
\textbf{Random Forest.}\quad
n\_estimators:\ 100,\ 200,\ 300,\ 400,\ 500;\quad
max\_depth:\ 8,\ 12,\ 16,\ 20,\ 24,\ 28;\quad
min\_samples\_split:\ 2,\ 5,\ 10;\quad
min\_samples\_leaf:\ 1,\ 2,\ 4;\quad
max\_features:\ sqrt,\ log2;\quad
class\_weight:\ balanced,\ balanced\_subsample.\\[-0.1em]
\textbf{XGBoost.}\quad
n\_estimators:\ 100,\ 200,\ 300,\ 400,\ 500;\quad
max\_depth:\ 4,\ 6,\ 8,\ 10,\ 12;\quad
learning\_rate:\ 0{.}01,\ 0{.}03,\ 0{.}05,\ 0{.}1;\quad
subsample:\ 0{.}5,\ 0{.}7,\ 0{.}85,\ 1{.}0;\quad
colsample\_bytree:\ 0{.}5,\ 0{.}7,\ 0{.}85,\ 1{.}0.
\end{latin}

%==============================================================================
\section{آزمایش \pnum{5.8}: گسترش‌های حذف تدریجی \lr{ETHOS}}
\label{sec:ethos-ablation}
%==============================================================================

\subsection{هدف}

گسترش مطالعۀ حذف تدریجی برای آزمون اینکه آیا آموزش \emph{نظارت‌شده} (بدون کم‌نمونه اپیزودیک) عملکرد \lr{ETHOS} را بهبود می‌دهد. گونۀ اصلی \texttt{\lr{no\_few\_shot}} (\lr{NN} استاندارد) به \lr{AUROC} \pnum{0.645} رسید؛ فرض می‌کنیم \lr{ETHOS} با سر طبقه‌بندی نظارت‌شده ممکن است از آموزش اپیزودیک پیشی گیرد.

\subsection{طراحی}

رمزگذار \lr{ETHOS} را با سر نظارت‌شده آموزش می‌دهیم: \lr{Linear(128$\to$64$\to$1)} با سیگموئید، \lr{CrossEntropyLoss} با وزن کلاس. همان تقسیم‌های داده که \lr{RF} (\lr{seed=42}). \pnum{100} دوره، صبر توقف زودهنگام \pnum{15}. اختیاری: بارگذاری رمزگذار \lr{ETHOS} پیش‌آموزش‌دیده از پیش‌آموزش \lr{MLM}.

\subsection{خلاصۀ نتایج}

\texttt{\lr{ETHOS\_Supervised}} به \lr{AUROC} \pnum{0.610}، \lr{F1} \pnum{0.386}، \lr{Accuracy} \pnum{0.627} می‌رسد. در مقایسه با \lr{Proposed\_FewShot} در \lr{Exp1} (\pnum{0.611})، آموزش نظارت‌شده \lr{AUROC} تقریباً هم‌تراز می‌دهد. سر نظارت‌شده از کل مجموعه‌داده بدون نمونه‌برداری اپیزودیک یاد می‌گیرد. جدول~\ref{tab:exp1} و جدول~\ref{tab:exp8} را ببینید.

%==============================================================================
\section{آزمایش \pnum{5.9}: معماری‌های ترکیبی}
\label{sec:exp-hybrid}
%==============================================================================

\subsection{هدف}

مقایسۀ معماری‌های ترکیبی که بردارهای نهفتۀ \lr{ETHOS} را با ویژگی‌های جدولی ترکیب می‌کنند. آزمون: (\pnum{1}) \lr{RF} روی [\lr{ETHOS} \pnum{128}ب + جدولی \pnum{25}ب] = \pnum{153}ب، (\pnum{2}) \lr{RF} تنها روی \lr{ETHOS} \pnum{128}ب.

\subsection{طراحی}

\begin{enumerate}
\item \textbf{معماری ترکیبی \lr{ETHOS} و \lr{Random Forest}}: استخراج بردارهای نهفتۀ \lr{ETHOS} برای همه‌ی نمونه‌ها. الحاق با ویژگی‌های جدولی: $\mathbf{x}_{\text{\lr{hybrid}}} = [\mathbf{e}_{\text{\lr{ETHOS}}}; \mathbf{x}_{\text{\lr{tab}}}] \in \mathbb{R}^{153}$. آموزش \lr{Random Forest} با \lr{RandomizedSearchCV} (\pnum{50} آزمون) روی آموزش+اعتبارسنجی. تنظیم آستانه روی اعتبارسنجی.
\item \textbf{\lr{Random Forest} فقط روی نهفته‌سازی \lr{ETHOS}}: آموزش \lr{Random Forest} تنها روی بردارهای نهفتۀ \pnum{128}-بعدی \lr{ETHOS} (بدون جدولی). همان \lr{CV} و تنظیم آستانه.
\end{enumerate}

\subsection{خلاصۀ نتایج}

معماری ترکیبی \lr{ETHOS} و \lr{Random Forest} به \lr{AUROC} \pnum{0.717}، \lr{F1} \pnum{0.652}، \lr{Accuracy} \pnum{0.665} می‌رسد؛ بهبود چشمگیر نسبت به \lr{Proposed\_FewShot} در \lr{Exp1} (\pnum{0.611}) و \texttt{\lr{ETHOS\_Supervised}} (\pnum{0.610}). ترکیبی هم بازنمایی‌های آموخته‌شده و هم ویژگی‌های مهندسی‌شده را بهره می‌برد و شکاف تا \lr{RF} (\pnum{0.760}) را به حدود \pnum{0.04} \lr{AUROC} کاهش می‌دهد. \lr{Random Forest} فقط روی نهفته‌سازی \lr{ETHOS} به \lr{AUROC} \pnum{0.598} و \lr{F1} \pnum{0.576} می‌رسد؛ ضعیف‌تر از ترکیبی است و تأیید می‌کند ویژگی‌های جدولی سیگنال بحرانی اضافه می‌کنند. جدول~\ref{tab:exp8} را ببینید.

%==============================================================================
\section{جمع‌بندی فصل}
%==============================================================================

این فصل چیدمان تجربی را شرح داد: پیاده‌سازی و معماری نرم‌افزار، پیکربندی داده، پیکربندی چارچوب \lr{ETHOS} و یادگیری کم‌نمونه و خط پایۀ کلاسیک، فرضیات آزمایش، توان آماری و اندازۀ نمونه، مشخصات سخت‌افزار و نرم‌افزار، فلسفۀ طراحی، رویۀ اجرا، پیکربندی‌های مدل خط پایه، معیارهای ارزیابی و گزارش، و طراحی تفصیلی برای هر یک از آزمایش‌های \pnum{1} تا \pnum{5}، گسترش‌های حذف تدریجی \lr{ETHOS}، و معماری‌های ترکیبی. فصل بعد نتایج را گزارش می‌کند.
